\documentclass[11pt]{article}

\usepackage[T1]{fontenc}
\usepackage[utf8]{inputenc}
\usepackage{lmodern}
\usepackage[a4paper,margin=27mm]{geometry}
\usepackage{microtype}
\usepackage{amsmath,amssymb,mathtools}
\usepackage{enumitem}
\usepackage{booktabs,tabularx}
\usepackage[colorlinks=true,linkcolor=blue,citecolor=blue,urlcolor=blue]{hyperref}

\title{Projection-First Dark-Sector Diagnostics:\\
\large Effective Stress, Lensing, and Cosmological Closure}
\author{Peter Nero}
\date{July 2026\\Version 2}

\begin{document}
\maketitle

\begin{abstract}
Dark matter and dark energy are names for distinct empirical requirements,
not merely for missing microscopic descriptions. A projection-first
framework may suggest that unresolved source data appear as effective
gravitational terms, but this suggestion becomes physical only after it
produces covariant field equations and observables. This paper formulates
that requirement for Modal Triplet Theory. A dark-matter realization must
reproduce gravitational lensing, approximately pressureless clustering,
structure growth, cosmic microwave background effects, and cluster
dynamics. A dark-energy realization must produce the observed expansion
history and perturbation behavior while remaining stable and compatible
with local gravity. Both require a selected effective stress tensor or an
explicit modified-gravity operator satisfying the Bianchi identity.
``Constraint load'' and ``capacity exhaustion'' are therefore retained as
hypothesis labels, not established explanations. No current MTT result
derives the required dark-sector source, halo solutions, or cosmological
likelihoods.
\end{abstract}

\section*{Version 2 Revision Note}
\begin{description}[leftmargin=3.2cm,style=nextline]
\item[Supersedes] Version 1.0, \emph{A Projection-First Reframing of Dark
Matter and Dark Energy}.
\item[Reason] The original paper identified failed particle encoding with
dark matter and global capacity exhaustion with dark energy without deriving
the stress-energy, field equations, perturbations, or observables needed for
those identifications.
\item[Resolution] The two proposals are recast as falsifiable hypotheses.
Necessary covariant source equations, cold-matter limits, lensing tests,
acceleration conditions, stability gates, and data comparisons are stated
explicitly.
\item[Retained result] Projection-first reasoning can motivate a search for
effective gravitational sources that need not be new elementary particles.
\item[Remaining boundary] MTT has not selected a dark-sector action or
effective tensor and has not reproduced halo, lensing, CMB, growth, BAO, or
supernova likelihoods.
\end{description}

\tableofcontents

\section{The empirical problem}

Dark matter and dark energy are inferred from different classes of
observations. Dark-matter phenomenology includes galaxy and cluster
dynamics, gravitational lensing, structure formation, and signatures in
the cosmic microwave background. Dark-energy phenomenology concerns the
late-time expansion history and the behavior of cosmological perturbations.
The spatial separation between lensing mass and collisional gas in merging
clusters is an especially important constraint on any account of gravitating
mass \cite{Clowe2006}. Precision cosmology is well described by a
\(\Lambda\)CDM baseline, although time-dependent dark-energy models remain
actively tested \cite{Planck2018,DESI2024}.

These observations do not require that dark matter be an elementary
particle or that dark energy be a literal fluid. They do require a theory to
produce the same gravitational observables. A semantic replacement of
``dark component'' by ``projection strain'' is not enough.

\section{The covariant closure requirement}

In an Einstein-form description one may write
\[
 G_{\mu\nu}+\Lambda g_{\mu\nu}
 =8\pi G\left(T^{\rm vis}_{\mu\nu}
              +T^{\rm eff}_{\mu\nu}\right).
\]
The Bianchi identity gives
\[
 \nabla^\mu G_{\mu\nu}=0.
\]
If visible matter is separately conserved, then consistency requires
\[
 \nabla^\mu T^{\rm eff}_{\mu\nu}=0.
\]
An exchange law is possible, but it must be specified:
\[
 \nabla^\mu T^{\rm vis}_{\mu\nu}=Q_\nu,
 \qquad
 \nabla^\mu T^{\rm eff}_{\mu\nu}=-Q_\nu.
\]

A modified-gravity description can instead use
\[
 \mathcal E_{\mu\nu}[g,\chi]
 =8\pi G T^{\rm vis}_{\mu\nu},
\]
where \(\chi\) denotes additional geometric or projected variables.
Moving terms between the two sides can define an effective tensor, but it
does not remove the need for a covariant identity, initial data, stability,
and observables.

For MTT, the missing central object is consequently one of:
\[
 S_{\rm selected}[g,\chi]
 \quad\Longrightarrow\quad
 T^{\rm eff}_{\mu\nu}
 =-\frac{2}{\sqrt{-g}}
   \frac{\delta S_{\rm selected}}{\delta g^{\mu\nu}},
\]
or a selected covariant operator \(\mathcal E_{\mu\nu}\) with an equivalent
closure identity. A scalar ``closure cost'' is not a substitute for this
tensor.

\section{Dark matter: required behavior}

\subsection{Background and perturbations}

At homogeneous level, cold dark matter is approximately pressureless:
\[
 p_{\rm dm}\simeq0,
 \qquad
 \rho_{\rm dm}\propto a^{-3}.
\]
At perturbative level, a proposed source must supply sufficiently small
effective sound speed and controlled anisotropic stress so that it clusters
on the observed scales. It must participate in the metric potentials that
govern both motion and lensing.

In a scalar-perturbed cosmology,
\[
 ds^2=-(1+2\Psi)dt^2
 +a(t)^2(1-2\Phi)d\mathbf x^2,
\]
nonrelativistic motion primarily probes \(\Psi\), while lensing probes
combinations of \(\Phi\) and \(\Psi\). Reproducing one rotation curve does
not establish the required lensing relation or cosmological growth law.

\subsection{Minimum dark-matter test set}

An MTT dark-matter candidate must provide:
\begin{enumerate}[leftmargin=*]
\item a covariant background solution and a positive effective density;
\item linear scalar, vector, and tensor perturbation equations;
\item galaxy and cluster lensing from the same metric that governs motion;
\item halo or alternative nonlinear solutions;
\item a matter power spectrum and CMB transfer functions;
\item merging-cluster behavior and bounds on dissipation or self-interaction;
\item compatibility with solar-system and binary-pulsar tests; and
\item a parameter ledger and likelihood comparison with \(\Lambda\)CDM.
\end{enumerate}

The phrase ``constraint load without particle encoding'' is compatible with
either an effective source or modified gravity. It does not imply that the
source is cold, collisionless, stable, or correctly distributed. Those are
outputs to be calculated.

\section{Dark energy: required behavior}

\subsection{Acceleration}

For a spatially homogeneous Einstein cosmology,
\[
 \frac{\ddot a}{a}
 =-\frac{4\pi G}{3}\left(\rho+3p\right).
\]
Acceleration requires
\[
 \rho+3p<0
\]
for the effective total source in this formulation. A cosmological constant
has \(p_\Lambda=-\rho_\Lambda\). More general dark energy is often described
by \(w(a)=p/\rho\), but a background \(w(a)\) alone is insufficient. One must
also specify perturbations, sound speed, anisotropic stress, and interactions.

The original ``capacity exhaustion'' language does not determine
\(\rho(a)\), \(p(a)\), or \(w(a)\), and therefore does not predict
acceleration. It remains a possible interpretation only after a selected
MTT action or closure equation yields these quantities.

\subsection{Minimum dark-energy test set}

A viable construction must supply:
\begin{enumerate}[leftmargin=*]
\item a background \(H(z)\) and luminosity and angular-diameter distances;
\item BAO, supernova, CMB, and growth predictions from one parameter set;
\item stable scalar and tensor perturbations without ghosts or gradient
instabilities;
\item a controlled effective Newton coupling and gravitational slip;
\item compatibility with local tests and gravitational-wave propagation;
\item an initial-condition and radiative-stability account; and
\item evidence that the new description improves prediction or parameter
economy.
\end{enumerate}
Current observations motivate testing both a cosmological constant and
dynamical alternatives; the paper does not assume that one option has
already been selected.

\section{What projection can and cannot contribute}

A projection
\[
 \pi:\mathcal X_{\rm upper}\longrightarrow\mathcal X_{\rm eff}
\]
can hide variables that influence the retained sector. Integrating out or
conditioning on hidden variables may yield effective forces, memory kernels,
noise, or additional source terms. This is a legitimate route to dark-sector
phenomenology.

However, the effective term depends on the upper action, state, projection,
and approximation. The same projection map can produce no force, a local
potential, colored noise, dissipation, or a nonlocal kernel under different
upper dynamics. Projection alone does not select the sign, equation of state,
clustering scale, or magnitude of a dark source.

The disciplined MTT claim is therefore:
\begin{quote}
Selected upper degrees of freedom may descend to an effective gravitational
source that is not represented as an elementary particle in the lower
description.
\end{quote}
Whether that source behaves as dark matter or dark energy is a subsequent
calculation.

\section{Two constructive routes}

\subsection{Effective-source route}

Start with a selected upper action
\[
 S[g,\psi,\chi],
\]
where \(\psi\) denotes visible fields and \(\chi\) denotes projected or
unresolved fields. Define a state and integrate or solve for \(\chi\) under a
controlled approximation:
\[
 e^{iS_{\rm eff}[g,\psi]}
 =\int\mathcal D\chi\;e^{iS[g,\psi,\chi]}.
\]
Then derive \(T^{\rm eff}_{\mu\nu}\), conservation laws, response kernels,
and perturbations. This route can yield particle, fluid, stochastic, or
nonlocal behavior.

\subsection{Modified-geometry route}

Derive a selected covariant operator \(\mathcal E_{\mu\nu}\) from the MTT
closure geometry. Prove its differential identity, count propagating
degrees of freedom, and establish a well-posed stable initial-value problem.
Then compute lensing and cosmology. Writing the result as an effective
stress tensor is optional; matching observables is not.

Both routes require the same-source upper data. Importing a fitted halo
profile or a fitted \(w(a)\) would demonstrate compatibility but not derive
the dark sector.

\section{Status and falsifiability}

\begin{center}
\small
\begin{tabularx}{\textwidth}{@{}p{0.31\textwidth}p{0.18\textwidth}X@{}}
\toprule
Claim & Status & Missing object or test\\
\midrule
Projection can generate effective sources & Standard possibility &
Requires upper dynamics, state, and reduction.\\
MTT closure cost is a stress tensor & Not established & Metric variation of
a selected covariant action.\\
MTT derives dark matter & Open & Cold clustering, lensing, CMB, halos, and
likelihoods.\\
MTT derives dark energy & Open & \(H(z)\), perturbations, stability, and
likelihoods.\\
Dark matter need not be a new particle & Logically possible & Must still
match all gravitational and cosmological evidence.\\
Dark energy is capacity exhaustion & Interpretive hypothesis & No selected
equation of state or source amplitude yet.\\
One MTT source explains both sectors & Open & Common action and independent
predictions required.\\
\bottomrule
\end{tabularx}
\end{center}

The hypotheses are falsifiable. A candidate fails if it cannot reproduce
lensing and dynamics from one metric, violates stability or local gravity,
spoils the CMB or matter spectrum, or needs observational dark-sector
profiles inserted as unconstrained functions.

\section{Discussion}

The projection-first perspective contributes a useful warning: the absence
of a lower particle encoding does not prove the absence of upper physical
influence. Effective theories routinely contain forces and stresses from
variables that have been removed. This opens a legitimate nonparticle route.

The same warning applies in reverse. A variable missing from a description
does not automatically gravitate, and a bookkeeping deficit does not
automatically accelerate the universe. General covariance turns that gap
into a concrete demand: produce the tensor or modified operator, prove its
identity and stability, and confront the data.

This standard is productive for MTT. It identifies a single mathematical
frontier rather than two verbal explanations: derive the selected
metric-response functional of the unresolved sector. Its homogeneous part
can then be tested as dark energy; its inhomogeneous and clustering part can
be tested as dark matter.

\section{Conclusion}

MTT has not yet explained dark matter or dark energy. It has a possible
organizing hypothesis: unresolved selected source data may appear in a
lower description as effective gravitational response rather than as new
elementary particles.

The next result must be covariant and numerical. It should derive
\(T^{\rm eff}_{\mu\nu}\) or \(\mathcal E_{\mu\nu}\) from a selected MTT
action, obtain background and perturbation equations, and compare their
lensing, growth, CMB, BAO, and supernova predictions with existing models.
Until then, ``constraint load'' and ``capacity exhaustion'' name research
directions, not dark-sector solutions.

\begin{thebibliography}{9}

\bibitem{Riess1998}
A.~G. Riess et al.,
\emph{Observational Evidence from Supernovae for an Accelerating Universe
and a Cosmological Constant},
Astronomical Journal \textbf{116} (1998), 1009--1038.

\bibitem{Perlmutter1999}
S. Perlmutter et al.,
\emph{Measurements of \(\Omega\) and \(\Lambda\) from 42 High-Redshift
Supernovae},
Astrophysical Journal \textbf{517} (1999), 565--586.

\bibitem{Clowe2006}
D. Clowe et al.,
\emph{A Direct Empirical Proof of the Existence of Dark Matter},
Astrophysical Journal Letters \textbf{648} (2006), L109--L113.

\bibitem{Planck2018}
Planck Collaboration,
\emph{Planck 2018 Results. VI. Cosmological Parameters},
Astronomy \& Astrophysics \textbf{641} (2020), A6.

\bibitem{DESI2024}
DESI Collaboration,
\emph{DESI 2024 VI: Cosmological Constraints from the Measurements of
Baryon Acoustic Oscillations},
Journal of Cosmology and Astroparticle Physics \textbf{02} (2025), 021.

\bibitem{Weinberg1989}
S. Weinberg,
\emph{The Cosmological Constant Problem},
Reviews of Modern Physics \textbf{61} (1989), 1--23.

\end{thebibliography}

\end{document}
